% !TEX root = ../algebra_02.tex

\section{Циклические подпространства}

\begin{Definition}{Циклическое подпространство}{}
	Пусть $v \in V$ --- произвольный вектор. Циклическим подпространством $C_v$ называется минимальное $A$-инвариантное подпространство, содержащее вектор $v$. Эквивалентно, оно определяется как пересечение всех таких подпространств:
	\[
	C_v = \bigcap\{W \subseteq V \mid W \text{ — $A$-инвариантное подпространство и } v \in W\}
	\].
\end{Definition}

Для изучения структуры $C_v$ вводится понятие аннулятора.

\begin{Definition}{Аннулятор вектора}{}
	Множество $I_v = \{ f \in K[x] \mid f(A)v = 0 \}$ называется аннулятором вектора $v$. Нормированный многочлен минимальной степени из $I_v$, обозначаемый $M_{v,A}(x)$, называется минимальным аннулятором вектора $v$ относительно оператора $A$.
\end{Definition}

\begin{Proposition}{Базис циклического подпространства}{}
	Пусть $\deg M_{v,A} = d$. Тогда система векторов $v, Av, A^2v, \dots, A^{d-1}v$
	является базисом циклического подпространства $C_v$.
\end{Proposition}

\begin{proof}
	Линейная независимость этой системы напрямую следует из минимальности степени многочлена $M_{v,A}$: если бы векторы были линейно зависимы, существовал бы ненулевой многочлен степени меньше $d$, аннулирующий $v$, что противоречит определению $M_{v,A}$. Кроме того, из уравнения $M_{v,A}(A)v = 0$ можно выразить вектор $A^dv$ через векторы меньших степеней $A^kv$, откуда следует, что указанная система порождает все подпространство $C_v$ и образует его базис.
\end{proof}
